\documentclass[10pt,aspectratio=169,fontset=none]{ctexbeamer}
\usepackage[T1]{fontenc}

\title{Subliming 主题 Beamer 示例}
\author{杨曜堃}
\institute{地球科学与信息物理学院}
\date{February 19, 2025}

% 添加封面上的大学校标，这里使用的是白色中南大学校标，并设置透明度为 30%
% 如果你想使用校徽，可以在 beamerinnerthemeSUBLIMING.sty 文件中修改，找到对应的注释即可
% 我的意思是，四个字的那种叫校徽，圆圆的那个叫校标
\titlegraphic{\includegraphics[width=8cm]{assets/csu-logo.pdf}}

\usepackage{styles/SUBLIMINGmacros}
\usefolder{styles}
\usetheme[style=csublue, cover=facet]{SUBLIMING} % 在这里更换颜色主题风格
% 支持 csublue 和 nupurple 两种风格

\begin{document}

    \begin{frame}
        \titlepage
    \end{frame}

    \begin{frame}{介绍 Introduction}
        Subliming 主题是一个为中南大学设计的 Beamer 主题，也可以通过调整模板中的颜色、字体和背景以适应其他大学的风格。 
        \\ ~ \\
        这个主题是在 Northwestern University 的 Aaron Wolf 制作的 \href{https://github.com/aarondwolf/wildcatwildcat}{Wildcat} 主题的基础上修改设计的，更适合国内高校。同其一样，Subliming 主题也遵循 \href{http://www.gnu.org/licenses/}{GNU 通用公共许可证}，这意味着任何人都可以免费自由使用这个主题！
        
        \setblockstyle{metropolis}
        \begin{tbox}{default}
            你也可以自定义这个主题！我在各个包中增加了中文注释，以便你更好地理解这个主题的设计。如果有任何问题或建议，欢迎向我们的 \href{https://github.com/PourRevenir/Subliming}{GitHub 仓库}提交 \href{https://github.com/PourRevenir/Subliming/issues}{issue} 或 \href{https://github.com/PourRevenir/Subliming/pulls}{pull request}！
        \end{tbox}
    \end{frame}

    \begin{frame}{目录 Contents}
        \tableofcontents
    \end{frame}

    \section{1 选项}

    \begin{frame}{颜色 Colors}
        Subliming 主题允许切换主题颜色，并内置了几种风格。默认的风格为中南蓝，可通过修改如下选项来切换主题颜色：
        
        \begin{center}
            \texttt{\textcolor{primary}{\textbackslash usetheme}[style=颜色风格]\alert{\{SUBLIMING\}}}
        \end{center}
        
        目前内置的颜色风格包括：

       \begin{itemize}
            \item \textcolor[HTML]{186199}{中南蓝 csublue}
            \item \textcolor[RGB]{078, 042, 132}{西北紫 nupurple}
            \item \textcolor[HTML]{A00000}{湖大红 hnured}
        \end{itemize}  
       
    \end{frame}
    
    \begin{frame}{封面样式 Cover Style}
        你也可以自定义封面样式。默认的封面样式为 Facet，你可以通过修改如下选项来切换封面样式：
        \begin{center}
            \texttt{\textcolor{primary}{\textbackslash usetheme}[cover=封面样式]\alert{\{SUBLIMING\}}}
        \end{center}

        目前内置的封面样式包括低面数风格 facet 和空白 plain 两种，默认为 facet。
    \end{frame}


    \section{2 文字}

    \begin{frame}{字体和样式 Fonts \& Styles}
        内置了几种中英文文字样式，需要确保字体文件处于正确的路径下。
        \begin{columns}
            \column{0.3\textwidth}
            \begin{itemize}
                \item Regular
                \item \alert{Alert}
                \item \example{Example}
                \item \textit{Italic}
                \item \textbf{Bold}
                \item \textbf{\textit{Bold Italic}}
            \end{itemize}
            \column{0.6\textwidth}
            \begin{itemize}
                \item 常规
                \item \alert{醒目}
                \item \example{例子}
                \item \textit{斜体}
                \item \textbf{粗体}
                \item \textbf{\textit{粗斜体}}
            \end{itemize}
        \end{columns}
        \vspace*{0.5cm}
        参考文献引用：\cite{Mosher2013} \cite{Asgharzadeh2008} \cite{Frese1999} \cite{Gharti2017}
    \end{frame}

    \begin{frame}{数学字体 Math Fonts}
        使用了 \href{https://github.com/firamath/firamath}{Fira Math} 字体作为数学字体，这是一个开源的数学字体，适用于数学公式。需要在 Texlive 中安装， 可使用 \texttt{\alert{tlmgr} install firamath} 命令。

        \begin{equation}
            \int_{-\infty}^{\infty} e^{-x^2} dx = \sqrt{\pi}
        \end{equation}

        如果需要在 MixTex 中安装，则使用 \texttt{\alert{mpm} --admin --install=firamath} 命令。

    \end{frame}

    \begin{frame}{文本颜色 Text Color}
        你也可以自定义文本颜色。
        \\ ~ \\
        默认的文本颜色为黑色，你可以通过命令\texttt{\textbackslash textcolor\{颜色\}\{文本\}}得到彩色文本，内置颜色包括：
        \begin{columns}
            \column{0.3\textwidth}
            \begin{itemize}
                \item[$\textcolor{primary}{\bullet}$] \textcolor{primary}{primary} 
                \item[$\textcolor{primary40}{\bullet}$] \textcolor{primary40}{primary40}
                \item[$\textcolor{primary30}{\bullet}$] \textcolor{primary30}{primary30}
                \item[$\textcolor{primary20}{\bullet}$] \textcolor{primary20}{primary20}
                \item[$\textcolor{primary10}{\bullet}$] \textcolor{primary10}{primary10}
            \end{itemize}
            \column{0.3\textwidth}
            \begin{itemize}
                \item[$\textcolor{nubrightgreen}{\bullet}$] \textcolor{nubrightgreen}{nubrightgreen}
                \item[$\textcolor{nubrightteal}{\bullet}$] \textcolor{nubrightteal}{nubrightteal}
                \item[$\textcolor{nubrightblue}{\bullet}$] \textcolor{nubrightblue}{nubrightblue}
                \item[$\textcolor{nubrightyellow}{\bullet}$] \textcolor{nubrightyellow}{nubrightyellow}
                \item[$\textcolor{nubrightorange}{\bullet}$] \textcolor{nubrightorange}{nubrightorange}
                \item[$\textcolor{nubrightred}{\bullet}$] \textcolor{nubrightred}{nubrightred}
            \end{itemize}
            \column{0.3\textwidth}
            \begin{itemize}
                \item[$\textcolor{nudarkgreen}{\bullet}$] \textcolor{nudarkgreen}{nudarkgreen}
                \item[$\textcolor{nudarkteal}{\bullet}$] \textcolor{nudarkteal}{nudarkteal}
                \item[$\textcolor{nudarkblue}{\bullet}$] \textcolor{nudarkblue}{nudarkblue}
                \item[$\textcolor{nudarkyellow}{\bullet}$] \textcolor{nudarkyellow}{nudarkyellow}
                \item[$\textcolor{nudarkorange}{\bullet}$] \textcolor{nudarkorange}{nudarkorange}
                \item[$\textcolor{nudarkred}{\bullet}$] \textcolor{nudarkred}{nudarkred}
            \end{itemize}
        \end{columns}
    \end{frame}

    \section{3 盒子}

    \begin{frame}{tcolorbox 盒子示例}
        这里提供了三种彩色盒子，你可以通过 \texttt{\textcolor{primary}{\textbackslash begin}\{tblock\}} 环境创建. 你也可以使用 \texttt{talert} 和 \texttt{texample} 这两种环境分别生成其他颜色的盒子。
        \begin{tblock}{tblock}
            这里是一个 tblock 盒子，颜色将和主题色一致。
        \end{tblock}
        \begin{talert}{talert}
            你也可以动手修改，比如制作一个可以自动编号的定理环境。
        \end{talert}
        \begin{texample}{texample}
            这里是一个 texample 盒子。
        \end{texample}
    \end{frame}

    \begin{frame}[fragile]{可自定义颜色的 tcb 盒子 }
        使用\texttt{tfacetbox} 可以自定义盒子的颜色。 这只适用于非原色（红、绿、蓝），因为不容易进行阴影处理。
        \begin{tfacetbox}[nudarkyellow]{自定义颜色 Facet 盒子}
            这是一个 tfacetbox 盒子. 可通过以下代码创建:
            \begin{minted}{latex}
    \begin{tfacetbox}[自定义的颜色]{盒子标题}
        这是一个 tfacetbox 盒子...
    \end{tfacetbox}
            \end{minted}
        \end{tfacetbox}
    \end{frame}

    \begin{frame}{盒子样式(I) Box Styles}
        这里提供几种没有圆角和阴影的盒子样式，你可以通过 \texttt{\textcolor{primary}{\textbackslash setblockstyle}\{样式名称\}} 命令来切换盒子样式。

        \begin{center}
            \texttt{\textcolor{primary}{\textbackslash setblockstyle}\{native\}}
        \end{center}

        \setblockstyle{native}
        \centering
        \begin{minipage}[b]{0.5\textwidth}
            \begin{block}{Default}
                这是一个 Default 盒子
            \end{block}
            \begin{alertblock}{Alert}
                这是一个 Alert 盒子
            \end{alertblock}
            \begin{exampleblock}{Example}
                这是一个 Example 盒子
            \end{exampleblock} 
        \end{minipage}

    \end{frame}

    \begin{frame}{盒子样式(II) Box Styles}
    如果使用    

        \begin{center}
            \texttt{\textcolor{primary}{\textbackslash setblockstyle}\{nobackground\}}
        \end{center}
    则会得到无背景的盒子样式。
        \setblockstyle{nobackground}
        \centering
        \begin{minipage}[b]{0.5\textwidth}
            \begin{block}{Default}
                这是一个 Default 盒子
            \end{block}
            \begin{alertblock}{Alert}
                这是一个 Alert 盒子
            \end{alertblock}
            \begin{exampleblock}{Example}
                这是一个 Example 盒子
            \end{exampleblock}       
        \end{minipage}

    \end{frame}

    \begin{frame}{盒子样式(III) Box Styles}
    下面是 Metropolis 风格的盒子样式。    

        \begin{center}
            \texttt{\textcolor{primary}{\textbackslash setblockstyle}\{metropolis\}}
        \end{center}

        \setblockstyle{metropolis}
        \centering
        \begin{minipage}[b]{0.5\textwidth}
            \begin{block}{Default}
                这是一个 Default 盒子
            \end{block}
            \begin{alertblock}{Alert}
                这是一个 Alert 盒子
            \end{alertblock}
            \begin{exampleblock}{Example}
                这是一个 Example 盒子
            \end{exampleblock}       
        \end{minipage}

    \end{frame}

    \begin{frame}{盒子样式(IV) Box Styles}
    以及 Emph 风格的盒子样式。   

        \begin{center}
            \texttt{\textcolor{primary}{\textbackslash setblockstyle}\{emph\}}
        \end{center}

        \setblockstyle{emph}
        \centering
        \begin{minipage}[b]{0.5\textwidth}
            \begin{block}{Default}
                这是一个 Default 盒子
            \end{block}
            \begin{alertblock}{Alert}
                这是一个 Alert 盒子
            \end{alertblock}
            \begin{exampleblock}{Example}
                这是一个 Example 盒子
            \end{exampleblock}       
        \end{minipage}
    \end{frame}

    \begin{frame}[fragile]{tbox 盒子}
        你也可以使用 \texttt{tbox} 环境来创建文本盒子，作为一种无标题的盒子样式。按照如下方式使用
        
        \begin{minipage}[b]{0.23\textwidth}
        \setblockstyle{native}
        \begin{tbox}{default}
            Default 文本
        \end{tbox}
        \begin{tbox}{alerted}
            Alert 文本
        \end{tbox}
        \begin{tbox}{example}
            Example 文本
        \end{tbox}      
        \end{minipage}
        \hspace{1cm}
        \begin{minipage}[b]{0.23\textwidth}
        \setblockstyle{metropolis}
        \begin{tbox}{default}
            Default 文本
        \end{tbox}
        \begin{tbox}{alerted}
            Alert 文本
        \end{tbox}
        \begin{tbox}{example}
            Example 文本
        \end{tbox}      
        \end{minipage}
        \hspace{0.7cm}
        \begin{minipage}[b]{0.35\textwidth}
            \begin{minted}{latex}
\setblockstyle{盒子样式}
\begin{tbox}{盒子颜色}
    文本内容
\end{tbox}
            \end{minted}
        \end{minipage}
    \end{frame}

    \begin{frame}{盒子样式 (Default)}
        你也可以使用 Beamer 提供的默认文本盒子样式。
        \begin{center}
            \begin{minipage}[b]{0.5\textwidth}
                \begin{block}{ Default }
                    这里是一个 Default 盒子
                \end{block}
                \begin{alertblock}{Alert}
                    这里是一个 Alert 盒子
                \end{alertblock}
                \begin{exampleblock}{Example}
                    这里是一个 Example 盒子
                \end{exampleblock}
            \end{minipage}
        \end{center}
    \end{frame}

    \section{4 图表}

    \begin{frame}{表格 Tables}
            \begin{table}
            \caption{Largest cities in the world (source: Wikipedia)}
            \label{tab:table1}
            \begin{tabular}{@{} lr @{}}
            \toprule
            City & Population\\
            \midrule
            Mexico City & 20,116,842\\
            Shanghai & 19,210,000\\
            Peking & 15,796,450\\
            Istanbul & 14,160,467\\
            \bottomrule
            \end{tabular}
            \hspace*{1cm}
              \setlength\extrarowheight{3pt}
            \begin{tabular}{|lr|}
            \hline
            \rowcolor{primary}\color{background}City & \color{background}Population\\
            \hline
            Mexico City & 20,116,842\\
            Shanghai & 19,210,000\\
            Peking & 15,796,450\\
            Istanbul & 14,160,467\\
            \hline
            \end{tabular}
        \end{table}
    \end{frame}

    \begin{frame}{绘图 Plots}
       
    
    \end{frame}

    \section{5 参考文献}

    \begin{frame}[allowframebreaks,noframenumbering]{参考文献 References}
        \bibliographystyle{SUBLIMING} 
        \bibliography{references}
    \end{frame}

    \begin{frame}{标准页 Standard Page}
        你可以通过 \texttt{\textcolor{primary}{\textbackslash standout}\{文本\}} 命令创建一个标准页。文本内容将居中显示。
    \end{frame}

    \standout{Questions \& Answers}

\end{document}